\section{Motivation}

Earth-observation missions are moving from isolated spacecraft operations toward constellations in which sensing, communication, and decision-making are distributed across many nodes. This transition creates two coupled design problems. The first is geometric: a spacecraft must obtain access to a target at an operationally useful time. The second is informational: the task describing that target must reach a suitable spacecraft early enough for the access to be used. A constellation with excellent coverage can therefore fail a time-critical mission if the task is disseminated too late, while a well-connected network cannot observe an event for which no suitable access exists.

Wildfire response makes this coupling visible. A remote-sensing product may indicate a new or growing fire, but the detection alone is not the end of the operational chain. The event must be transformed into a task, prioritized, released into the network, forwarded within finite communication opportunities, received by one or more candidate observers, and finally executed before a response deadline. A useful performance evaluation must preserve this event order. It must also distinguish the time to first reception, the breadth of dissemination, and the completion of a follow-up observation.

One way to improve coordination is to assign selected spacecraft a central-node role. Conceptually, a central node may have additional processing, storage, cross-link capability, routing responsibility, coordination authority, or connectivity to ground. Such a node can accelerate information exchange, but it can also increase hardware differentiation and create operational dependencies. The design question is therefore not simply whether central nodes help. It is whether a particular degree of centralization improves mission fulfilment sufficiently to justify the added resource burden and the possible loss of robustness.

This thesis investigates that question through a full-factorial Walker-constellation campaign driven by real, de-duplicated Sentinel-3 wildfire tasks. The implemented experiment varies the fraction of satellites labelled as central nodes while also varying constellation size, plane count, altitude, and inclination. The result is a controlled architecture database in which communication, reception, observation, and resource metrics can be evaluated together.

\section{Problem Statement}

The starting task-allocation logic was developed for a different Earth-observation application. Adapting it to wildfire detection required more than changing the target coordinates. The wildfire use case introduces multiple simultaneous events, unequal priority, priority-dependent deadlines, priority-dependent task sizes, region-specific task density, and the need to validate tasks against physical fire observations. It also requires the simulator to distinguish dissemination from observation fulfilment.

The central problem addressed by this work is stated as follows:

\begin{quote}
Given a time-stamped set of prioritized wildfire tasks, a Walker observer constellation, a ground-station configuration, finite communication opportunities, and a configurable fraction of central nodes, determine how architecture and centralization affect task reception, dissemination breadth, follow-up observation, communication burden, and decision robustness.
\end{quote}

The phrase \emph{decision robustness} is used for stability under objective weights. Physical robustness refers separately to the retention of mission performance after injected node failures. Keeping those two meanings separate prevents a stable ranking from being mistaken for a failure-tolerant architecture.

\section{Aim and Objectives}

The aim is to establish a defensible performance--resource--robustness trade space for wildfire task allocation in a distributed satellite system. The objectives are:

\begin{enumerate}
\item construct a traceable wildfire-task dataset from Sentinel-3 SLSTR FRP products, supported by FIRMS screening and event validation;
\item adapt the task-allocation framework to multiple release times, priority classes, task sizes, and deadlines;
\item model finite, time-dependent communication windows and evaluate observation only after task reception;
\item execute and validate the complete 891-case Walker and central-node design grid for two event regions and the expanded 539-scenario Fresh V2 campaign where available;
\item quantify central-node, orbit, constellation-size, communication-load, and regional effects without assigning causal meaning to unadjusted correlations;
\item identify exact Pareto sets and evaluate the stability of architecture rankings under alternative objective weights; and
\item quantify robustness through explicit random and targeted node-removal, link-window, ground-station, and capacity-degradation experiments rather than through nominal failure labels.
\end{enumerate}

Table~\ref{tab:objective_evidence} maps each objective to the evidence used in this draft. The two-region nominal layer is complete. Fresh V2 robustness, task-size, policy, and revised weight analyses are complete for California; the corresponding India Fresh V2 extension is not locally available and is not inferred.

\begin{table}[htbp]
\centering
\small
\caption{Objectives, evidence products, and status in this working draft.}
\label{tab:objective_evidence}
\begin{tabularx}{\textwidth}{L{0.07\textwidth}Y L{0.28\textwidth} L{0.17\textwidth}}
\toprule
No. & Objective & Primary evidence & Status \\
\midrule
O1 & Sentinel-3 wildfire task construction & De-duplicated 805-task CSV and processing history & Complete \\
O2 & Multi-task priority and deadline model & Case-level \path{wildfire_tasks_used.csv} & Complete \\
O3 & Capacity-aware dissemination and observation-after-reception & Windows, transfers, receptions, and observation results & Complete \\
O4 & Full architecture campaign & 1,782 earlier nominal cases; California Fresh V2 scenario archive & Nominal complete; Fresh V2 California complete \\
O5 & Architecture and regional effects & Earlier paired-region tables and Fresh V2 California factor effects & Partial for final V2 regional comparison \\
O6 & Pareto and weight sensitivity & Exact Pareto membership and 10,000 revised weight draws & Complete for Fresh V2 California \\
O7 & Failure-injection robustness & Random/targeted failures, outages, and capacity degradation & Complete for Fresh V2 California \\
\bottomrule
\end{tabularx}
\end{table}

\section{Research Questions}

The approved questions are made operational through the exported metrics:

\begin{enumerate}[label=RQ\arabic*.,leftmargin=*]
\item How can a PBL-focused distributed task-allocation framework be adapted to generate and process multiple time-critical wildfire tasks from real remote-sensing detections?
\item How do real wildfire task distributions, confidence values, fire intensity, and regional task density affect dissemination and observation fulfilment?
\item How does the fraction of central nodes influence first-reception latency, full dissemination, deadline success, observation fulfilment, routing depth, communication load, and architecture dependency?
\item Which Walker constellation configurations provide the best trade-off among performance, complexity, and robustness for wildfire monitoring?
\item How stable is the preferred central-node fraction under changes in task size, observation radius, failure mode, and objective-function weights?
\end{enumerate}

The available evidence answers RQ1 and the California-specific portions of RQ2--RQ5. Fresh V2 supplies validated California policy, task-size, failure, and weight experiments. A final Fresh V2 answer about regional load/generalization requires the India archive; observation-radius sensitivity and a separately extracted first-reception latency also remain open.

\section{Hypotheses}

Five hypotheses organize the analysis:

\begin{description}[leftmargin=3.5cm,style=nextline]
\item[H1: Diminishing returns] Increasing central-node fraction initially reduces dissemination latency and improves task fulfilment, but the marginal benefit diminishes after the network becomes sufficiently connected.
\item[H2: Resource penalty] Central-node count, communication opportunity, and dependency increase with centralization, so the fastest architecture need not be the preferred architecture.
\item[H3: Load dependence] A dense task region produces different capacity and ranking behaviour from a sparse event region even when the architecture grid is identical.
\item[H4: Geometry dependence] Successful reception is necessary but not sufficient for observation because a receiving satellite still needs a valid access before the deadline.
\item[H5: No universal optimum] The preferred case changes with region and objective weights; a defensible result is a Pareto set or stable candidate set rather than a universal scalar optimum.
\end{description}

\section{Scope and Boundaries}

Table~\ref{tab:scope_hierarchy_new} defines the scope hierarchy. Central-node fraction is the primary variable. Walker geometry supports the centrality analysis, while task size, observation radius, and node failure test the sensitivity of the main conclusion. A separate heterogeneous relay layer is outside the implemented nominal campaign.

\begin{table}[htbp]
\centering
\caption{Hierarchy of the implemented and planned thesis scope.}
\label{tab:scope_hierarchy_new}
\begin{tabularx}{\textwidth}{L{0.17\textwidth}L{0.31\textwidth}Y}
\toprule
Level & Study element & Treatment in this report \\
\midrule
Core & Central-node fraction in a wildfire DSS & Complete 0--100\% sweep for every base architecture \\
Supporting & Satellite count, plane count, altitude, inclination & Full-factorial nominal screening \\
Sensitivity & Objective weights & Complete 10,000-draw Fresh V2 California analysis \\
Sensitivity & Task size and injected failures & Complete for Fresh V2 California \\
Sensitivity & Observation radius & Fixed at 700 km; targeted sweep not executed \\
Extension & Dedicated heterogeneous central-node orbit & Conceptual discussion and future work only \\
Out of scope & Onboard fire-detection latency, weather, cloud, pointing, energy, and monetary cost & Explicitly excluded from final claims \\
\bottomrule
\end{tabularx}
\end{table}

The simulated mission chain starts when a processed wildfire task is injected. It does not simulate the original thermal measurement, onboard detection, cloud screening, product generation latency, or delivery of the final alert to civil-protection authorities. The observation radius is a mission-level access approximation and not a detailed optical-instrument pointing model. ``Cost'' means a declared non-monetary resource proxy; no procurement or lifecycle cost is inferred.

\section{Meaning of a Central Node}

The conceptual system definition and the implemented experiment are intentionally separated. Conceptually, central nodes may differ from ordinary satellites in hardware, processing, communication links, routing, coordination, and connection to ground. The final simulator represents only the portion for which validated code and outputs exist: selected nodes receive a central-node label, communication opportunities follow the configured link policy, and the routing policy preferentially forwards through the central-node structure. The campaign does not assign different sensor performance, processor speed, storage volume, reliability, or monetary cost to central nodes.

This distinction is essential for interpretation. A measured central-node effect is an effect of the implemented communication and routing role under common spacecraft geometry. It is not evidence that a particular hardware implementation is optimal. Hardware and operations can later be attached to the architectural role once real subsystem assumptions are available.

\section{Scientific and Engineering Contributions}

The thesis contributions are:

\begin{enumerate}
\item a real-data, de-duplicated Sentinel-3 task set that connects FRP, confidence, spatial clustering, priority, data size, and response deadline;
\item a modular DSS workflow separating physical communication and observation opportunities from the allocation/evaluation layer;
\item a capacity-aware and time-ordered forwarding evaluation that records task reception, transfer, observation, and failure evidence;
\item a complete two-region, 891-case-per-region nominal campaign plus a California Fresh V2 campaign with operational, policy, task-size, failure, and combined-stress scenarios;
\item a clear distinction between observation fulfilment, conditional reception latency, receiver coverage, and useful-recipient full dissemination;
\item exact Pareto filtering and Monte Carlo objective-weight sensitivity without claiming a universal optimum; and
\item reproducible random and targeted failure experiments with recorded seeds, separate original/surviving denominators, and time-aware mission metrics.
\end{enumerate}

The author developed the wildfire pipeline, integration, simulation automation, campaign execution, post-processing, statistical analysis, and report. Vincenzo Messina provided the communication model and the TAT-C basis. This allocation is stated here because the software components directly shape the interpretation of the results.

\section{Report Structure}

Chapter~2 presents the literature review. Chapter~3 explains the Sentinel-3 task dataset and de-duplication. Chapter~4 formalizes the Fresh V2 architecture, communication, routing, observation, and metrics. Chapter~5 separates the earlier two-region nominal evidence from the expanded Fresh V2 scenario design and validation. Chapter~6 reports the earlier validated regional baseline. Chapter~7 reports the available Fresh V2 California operational, policy, task-size, robustness, Pareto, and weight results. Chapter~8 answers the research questions with explicit evidence boundaries, while Chapter~9 gives the conclusions and remaining work. The appendices provide metric definitions, campaign configuration, validation records, supplementary tables, and software-release requirements.
